% !TEX TS-program = xelatex
% !BIB program = bibtex
% !TeX spellcheck = ru_RU
% !TEX root = talk.tex

\documentclass
  [ russian
  , aspectratio=1610 % Для защит онлайн лучше использовать разрешение не 4х3
  ] {beamer}

% !TeX spellcheck = ru_RU
% !TEX root = vkr.tex
% Опциональные добавления используемых пакетов. Вполне может быть, что они вам не понадобятся, но в шаблоне приведены примеры их использования.
\usepackage{tikz} % Мощный пакет для создание рисунков, однако может очень сильно замедлять компиляцию
\usetikzlibrary{decorations.pathreplacing,calc,shapes,positioning,tikzmark}

% Библиотека для TikZ, которая генерирует отдельные файлы для каждого рисунка
% Позволяет ускорить компиляцию, однако имеет свои ограничения
% Например, ломает пример выделения кода в листинге из шаблона
% \usetikzlibrary{external}
% \tikzexternalize[prefix=figures/]

\newcounter{tmkcount}

\usepackage{fancyvrb} % для verbatim переноса строки

% --- переносы в коде/JSON ---
\usepackage{listings}
\lstset{
  basicstyle=\ttfamily\small,
  breaklines=true,           % перенос длинных строк
  breakatwhitespace=false,   % переносить можно в любом месте
  columns=fullflexible,      % корректная ширина моношрифта
  keepspaces=true,           % сохранять пробелы
  showstringspaces=false
}
% (необязательно) Определение простого языка JSON для подсветки
\lstdefinelanguage{json}{
  string=[s]{"}{"},
  morestring=[s]{'}{'}, % если вдруг одинарные кавычки
  comment=[l]{//},
}

\tikzset{
    use tikzmark/.style={
            remember picture,
            overlay,
            execute at end picture={
                    \stepcounter{tmkcount}
                },
        },
    tikzmark suffix={-\thetmkcount}
}

\usepackage{booktabs} % Пакет для верстки "более книжных" таблиц, вполне годится для оформления результатов
% В шаблоне есть команда \multirowcell, которой нужен этот пакет.
\usepackage{multirow}
\usepackage{siunitx} % для таблиц с единицами измерений

% Для названий стоит использовать \textsc{}
\newcommand{\OCaml}{\textsc{OCaml}}
\newcommand{\miniKanren}{\textsc{miniKanren}}
\newcommand{\BibTeX}{\textsc{BibTeX}}
\newcommand{\vsharp}{\textsc{V$\sharp$}}
\newcommand{\fsharp}{\textsc{F$\sharp$}}
\newcommand{\csharp}{\textsc{C\#}}
\newcommand{\GitHub}{\textsc{GitHub}}
\newcommand{\SMT}{\textsc{SMT}}

\definecolor{eclipseGreen}{RGB}{63,127,95}
% \lstdefinelanguage{ocaml}{
% keywords={@type, function, fun, let, in, match, with, when, class, type,
% nonrec, object, method, of, rec, repeat, until, while, not, do, done, as, val, inherit, and,
% new, module, sig, deriving, datatype, struct, if, then, else, open, private, virtual, include, success, failure,
% lazy, assert, true, false, end},
% sensitive=true,
% commentstyle=\small\itshape\ttfamily,
% keywordstyle=\ttfamily\bfseries, %\underbar,
% identifierstyle=\ttfamily,
% basewidth={0.5em,0.5em},
% columns=fixed,
% fontadjust=true,
% literate={->}{{$\to$}}3 {===}{{$\equiv$}}1 {=/=}{{$\not\equiv$}}1 {|>}{{$\triangleright$}}3 {\\/}{{$\vee$}}2 {/\\}{{$\wedge$}}2 {>=}{{$\ge$}}1 {<=}{{$\le$}} 1,
% morecomment=[s]{(*}{*)}
% }

\makeatletter
\@ifclassloaded{beamer}{
    %%% Обязательные пакеты
    %% Beamer
    \usepackage{beamerthemesplit}
    \usetheme{SPbGU}
    \beamertemplatenavigationsymbolsempty
    \usepackage{appendixnumberbeamer}

    %% Локализация
    \usepackage{fontspec}
    \setmainfont{CMU Serif}
    \setsansfont{CMU Sans Serif}
    \setmonofont{CMU Typewriter Text}
    %\setmonofont{Fira Code}[Contextuals=Alternate,Scale=0.9]
    %\setmonofont{Inconsolata}
    \usepackage{polyglossia}
    \setmainlanguage{russian}
    \setotherlanguage{english}

    %% Графика
    \usepackage{pdfpages} % Позволяет вставлять многостраничные pdf документы в текст

    % Математические окружения с русским названием
    \newtheorem{rutheorem}{Теорема}
    \newtheorem{ruproof}{Доказательство}
    \newtheorem{rudefinition}{Определение}
    \newtheorem{rulemma}{Лемма}
    \usepackage{fancyvrb}
}
{}
\makeatother

\usepackage[autostyle]{csquotes} % Правильные кавычки в зависимости от языка
\usepackage{totcount}
\usepackage{setspace}
\usepackage{amsmath, amsfonts, amssymb, amsthm, mathtools} % "Адекватная" работа с математикой в LaTeX


\makeatletter

%!TEX root = vkr.tex

%% Параметры заполнения титульного листа
\usepackage{xkeyval}

%% Русскоязычный вариант
\define@key[ru]{mytitle}{chair}{\def\my@title@chair@ru{#1}}
\define@key[ru]{mytitle}{title}{\def\my@title@title@ru{#1}}
\define@key[ru]{mytitle}{group}{\def\my@title@group@ru{#1}}
\define@key[ru]{mytitle}{author}{\def\my@title@author@ru{#1}}
\define@key[ru]{mytitle}{supervisor}{\def\my@title@supervisor@ru{#1}}
\define@key[ru]{mytitle}{supervisorPosition}{\def\my@title@supervisorPosition@ru{#1}}
\define@key[ru]{mytitle}{reviewer}{\def\my@title@reviewer@ru{#1}}
\define@key[ru]{mytitle}{reviewerPosition}{\def\my@title@reviewerPosition@ru{#1}}
\define@key[ru]{mytitle}{consultant}{\def\my@title@consultant@ru{#1}}
\define@key[ru]{mytitle}{consultantPosition}{\def\my@title@consultantPosition@ru{#1}}
\define@key[ru]{mytitle}{year}{\def\my@title@year@ru{#1}}
\define@key[ru]{mytitle}{specialty}{\def\my@title@specialty@ru{#1}}
\define@key[ru]{mytitle}{programme}{\def\my@title@programme@ru{#1}}
\define@key[ru]{mytitle}{profile}{\def\my@title@profile@ru{#1}}
\define@choicekey*[ru]{mytitle}{type}{coursework,practice,prediploma,master,bachelor,production}{\def\my@title@type@ru{#1}}
\define@choicekey*[ru]{mytitle}{kind}{solution,experiment,production,comparison,theoretical}{\def\my@title@kind@ru{#1}}

%% Англоязычный вариант
\define@key[en]{mytitle}{chair}{\def\my@title@chair@en{#1}}
\define@key[en]{mytitle}{title}{\def\my@title@title@en{#1}}
\define@key[en]{mytitle}{group}{\def\my@title@group@en{#1}}
\define@key[en]{mytitle}{author}{\def\my@title@author@en{#1}}
\define@key[en]{mytitle}{supervisor}{\def\my@title@supervisor@en{#1}}
\define@key[en]{mytitle}{supervisorPosition}{\def\my@title@supervisorPosition@en{#1}}
\define@key[en]{mytitle}{reviewer}{\def\my@title@reviewer@en{#1}}
\define@key[en]{mytitle}{reviewerPosition}{\def\my@title@reviewerPosition@en{#1}}
\define@key[en]{mytitle}{consultant}{\def\my@title@consultant@en{#1}}
\define@key[en]{mytitle}{consultantPosition}{\def\my@title@consultantPosition@en{#1}}
\define@key[en]{mytitle}{year}{\def\my@title@year@en{#1}}
\define@key[en]{mytitle}{specialty}{\def\my@title@specialty@en{#1}}
\define@key[en]{mytitle}{programme}{\def\my@title@programme@en{#1}}
\define@key[en]{mytitle}{profile}{\def\my@title@profile@en{#1}}
\define@choicekey*[en]{mytitle}{type}{coursework,practice,prediploma,master,bachelor}{\def\my@title@type@en{#1}}
\define@choicekey*[en]{mytitle}{kind}{solution,experiment,production,comparison,theoretical}{\def\my@title@kind@en{#1}}

\newcommand{\filltitle}[2]{
    %% Значения по умолчанию для обоих языков
    \ifthenelse{\equal{#1}{ru}}
    {
        \presetkeys[#1]{mytitle}{
            year = {\the\year},
            type = {practice},
            reviewer = {},
            consultant = {},
            profile = {}
        }{}
    }
    {}
    \ifthenelse{\equal{#1}{en}}
    {
        \presetkeys[#1]{mytitle}{
            year = {\the\year},
            type = {practice},
            reviewer = {},
            consultant = {},
            profile = {}
        }{}
    }
    {}
    \setkeys[#1]{mytitle}{#2}
}


% !TeX spellcheck = ru_RU
% !TEX root = vkr.tex

%% Если что-то забыли, при компиляции будут ошибки Undefined control sequence \my@title@<что забыли>@ru
%% Если англоязычная титульная страница не нужна, то ее можно просто удалить.
\filltitle{ru}{
    %% Актуально только для курсовых/практик. ВКР защищаются не на кафедре а в ГЭК по направлению,
    %%   и к моменту защиты вы будете уже не в группе.
    chair              = {Кафедра информатики},
    group              = {23.Б16-мм},
    %
    %% Макрос filltitle ненавидит пустые строки, поэтому обязателен хотя бы символ комментария на строке
    %% Актуально всем.
    title              = {Реализация и внедрение архитектуры Gated Recurrent Unit в пакет DEGANN},
    %
    %% Здесь указывается тип работы. Возможные значения:
    %%   production - производственная практика;
    %%   coursework - отчёт по курсовой работе (ОБРАТИТЕ ВНИМАНИЕ, у техпрога и ПИ нет курсовых, только практики);
    %%   practice - отчёт по учебной практике;
    %%   prediploma - отчёт по преддипломной практике;
    %%   master - ВКР магистра;
    %%   bachelor - ВКР бакалавра.
    type               = {practice},
    %
    %% Здесь указывается вид работы. От вида работы зависят критерии оценивания.
    %%   solution - «Решение». Обучающемуся поручили найти способ решения проблемы в области разработки программного обеспечения или теоретической информатики с учётом набора ограничений.
    %%   experiment - «Эксперимент». Обучающемуся поручили изучить возможности, достоинства и недостатки новой технологии, платформы, языка и т. д. на примере какой-то задачи.
    %%   production - «Производственное задание». Автору поручили реализовать потенциально полезное программное обеспечение.
    %%   comparison - «Сравнение». Обучающемуся поручили сравнить несколько существующих продуктов и/или подходов.
    %%   theoretical - «Теоретическое исследование». Автору поручили доказать какое-то утверждение, исследовать свойства алгоритма и т.п., при этом не требуя написания кода.
    kind               = {production},
    %
    author             = {Мурадян Денис Степанович},
    %
    %% Актуально только для ВКР. Указывается код и название направления подготовки. Типичные примеры:
    %%   02.03.03 \enquote{Математическое обеспечение и администрирование информационных систем}
    %%   02.04.03 \enquote{Математическое обеспечение и администрирование информационных систем}
    %%   09.03.04 \enquote{Программная инженерия}
    %%   09.04.04 \enquote{Программная инженерия}
    %% Те, что с 03 в середине --- бакалавриат, с 04 --- магистратура.
    specialty          = {02.03.03 \enquote{Математическое обеспечение и администрирование информационных систем}},
    %
    %% Актуально только для ВКР. Указывается шифр и название образовательной программы. Типичные примеры:
    %%   СВ.5162.2020 \enquote{Технологии программирования}
    %%   СВ.5080.2020 \enquote{Программная инженерия}
    %%   ВМ.5665.2022 \enquote{Математическое обеспечение и администрирование информационных систем}
    %%   ВМ.5666.2022 \enquote{Программная инженерия}
    %% Шифр и название программы можно посмотреть в учебном плане, по которому вы учитесь.
    %% СВ.* --- бакалавриат, ВМ.* --- магистратура. В конце --- год поступления (не обязательно ваш, если вы были в академе/вылетали).
    programme          = {СВ.5162.2020 \enquote{Технологии программирования}},
    %
    %% Актуально всем.
    %% Должно умещаться в одну строчку, допускается использование сокращений, но без переусердствования,
    %% короткая строка с большим количеством сокращений выглядит странно
    %supervisorPosition = {проф. кафeдры системного программирования, д.ф.-м.н.,}, % Терехов А. Н.
    %supervisorPosition = {ст. преподаватель кафедры ИАС, к.~ф.-м.~н. (если есть),}, % Смирнов К. К.
    supervisorPosition = {доцент кафедры системного программирования, к.ф.-м.н.,},
    supervisor         = {Гориховский~В.~И.},
    %
    %% Актуально только для практик и курсовых. Если консультанта нет или он совпадает с научником, закомментировать или удалить вовсе.
    consultantPosition = {Преподаватель СПбГУ,},
    consultant         = {Алимов~П.~Г.},
}

% Английский титульник нужен только для ВКР, остальные виды работ могут его смело игнорировать.
\filltitle{en}{
    chair              = {Advisor's chair},
    group              = {ХХ.BХХ-mm},
    title              = {Template for SPbU qualification works},
    type               = {bachelor},
    author             = {FirstName Surname},
    %
    %% Possible choices:
    %%   02.03.03 \foreignquote{english}{Software and Administration of Information Systems}
    %%   02.04.03 \foreignquote{english}{Software and Administration of Information Systems}
    %%   09.03.04 \foreignquote{english}{Software Engineering}
    %%   09.04.04 \foreignquote{english}{Software Engineering}
    %% Те, что с 03 в середине --- бакалавриат, с 04 --- магистратура.
    specialty          = {02.03.03 \foreignquote{english}{Software and Administration of Information Systems}},
    %
    %% Possible choices:
    %%   СВ.5162.2020 \foreignquote{english}{Programming Technologies}
    %%   СВ.5080.2020 \foreignquote{english}{Software Engineering}
    %%   ВМ.5665.2022 \foreignquote{english}{Software and Administration of Information Systems}
    %%   ВМ.5666.2022 \foreignquote{english}{Software Engineering}
    programme          = {СВ.5162.2020 \foreignquote{english}{Programming Technologies}},
    %
    %% Note that common title translations are:
    %%   кандидат наук --- C.Sc. (NOT Ph.D.)
    %%   доктор ... наук --- Sc.D.
    %%   доцент --- docent (NOT assistant/associate prof.)
    %%   профессор --- prof.
    supervisorPosition = {Sc.D, prof.},
    supervisor         = {S.S. Supervisor},
    %
    consultantPosition = {position at \foreignquote{english}{Company}, degree if present},
    consultant         = {C.C. Consultant},
    %
    reviewerPosition   = {position at \foreignquote{english}{Company}, degree if present},
    reviewer           = {R.R. Reviewer},
}


% Определяем команды для титульного слайда
\newcommand{\academicGroup}{\my@title@group@ru}
\newcommand{\advisorChair}{\my@title@chair@ru}

\title[GRU в DEGANN]{Реализация и внедрение архитектуры GRU в пакет нейросетевой аппроксимации ДУ DEGANN}

% Авторы: ФИО с указанием СПбГУ
\author[Мурадян Д. С., Алимов П. Г.]{Мурадян Денис Степанович, Алимов Павел Геннадьевич}

% Дополнительные команды для титульного слайда
\newcommand{\publicationRecommendation}{Рекомендация к публикации: доцент Гориховский В. И.}

\makeatother
\date{}
\begin{document}
{
\setbeamertemplate{footline}{}
% Титульный слайд
\begin{frame}
    \includegraphics[width=1.4cm]{figures/SPbGU_Logo.png}
    \vspace{-35pt}
    \hspace{-10pt}
    \begin{center}
        \begin{tabular}{c}
            \scriptsize{Санкт-Петербургский государственный университет} \\[2pt]
            \scriptsize{\advisorChair}
        \end{tabular}
        \titlepage
    \end{center}

    \btVFill

    {\scriptsize
        {\publicationRecommendation}
    }
    \makeatother
    \begin{center}
        \vspace{5pt}
        \scriptsize{Санкт-Петербург\\ 2025}
    \end{center}
\end{frame}
}

% Слайд 1: Введение
\begin{frame}{Введение}
    \begin{itemize}
        \item Дифференциальные уравнения (ДУ) играют важную роль в науке и технике, однако численные методы их решения могут быть ресурсоёмкими и давать неточные результаты.
        \item Использование нейронных сетей для аппроксимации функций, описывающих ДУ, предлагает альтернативный подход.
        \item Функциональность DEGANN ограничена применением многослойного перцептрона (MLP), который не достаточно хорошо улавливает последовательные зависимости в данных.
        \item Рекуррентные нейронные сети, специально предназначены для работы с последовательностями и могут повысить эффективность аппроксимации.
    \end{itemize}
\end{frame}

% Слайд 2: Реализация и интеграция в DEGANN
\begin{frame}{Реализация и интеграция в DEGANN}
    \begin{itemize}
        \item Разработан класс \texttt{TensorflowGRUNet} на базе \texttt{tf.keras.Model}.
        \item Основные этапы реализации:
              \begin{itemize}
                  \item Создание списка GRU-слоёв с активацией \texttt{tanh} для основного состояния и \texttt{sigmoid} для рекуррентного.
                  \item Последовательная обработка входных данных через GRU-слои с последующим выводом через полносвязный слой с линейной активацией для регрессионного предсказания.
              \end{itemize}
        \item В модели реализованы методы для настройки оптимизатора, функции потерь, метрик и экспорта конфигурации, а также прямой проход данных.
        \item Интеграция в DEGANN осуществляется через параметр \texttt{net\_type} с сохранением общего интерфейса.
    \end{itemize}
\end{frame}

% Слайд 3: Валидация решения: Экспериментальная установка
\begin{frame}{Валидация решения: Экспериментальная установка}
    \begin{itemize}
        \item \textbf{Цель эксперимента:} показать, что при идентичных гиперпараметрах GRUNet обеспечивает более точное приближение целевых значений функций за счёт учета последовательных зависимостей.
        \item Сравнивались две архитектуры:
              \begin{itemize}
                  \item GRUNet (на основе GRU)
                  \item DenseNet (традиционный MLP)
              \end{itemize}
        \item Фиксированные параметры эксперимента:
              \begin{itemize}
                  \item 3 слоя, 30 нейронов в каждом.
                  \item Функция потерь: MSE; метрики: MAPE, \(R^2\).
                  \item Время обучения: 15, 30 и 45 секунд.
              \end{itemize}
    \end{itemize}
\end{frame}

% Слайд 4: Выбор тестовых функций и их обоснование
\begin{frame}{Выбор тестовых функций и их обоснование}
    \begin{itemize}
        \item Для проверки моделей выбраны сложные для аппроксимации функции, каждая из которых предъявляет уникальные требования:
              \begin{itemize}
                  \item \(\displaystyle f_{\sin}(x)=\sin(10x)\) – демонстрирует периодическую нелинейность, позволяющую оценить способность модели захватывать колебательные процессы.
                  \item \(\displaystyle f_{\text{hyperbol}}(x)=\frac{x^2+0.5}{x+0.1}\) – обладает неоднородной зависимостью, сочетая квадратичный рост с особенностями при малых значениях \(x\).
                  \item \(\displaystyle f_{\text{hardsin}}(x)=\sin\Bigl(\ln\bigl(x^{\sin(10x)}\bigr)\Bigr)\) – включает логарифмическое преобразование и экспоненциальные эффекты, что усиливает нелинейность.
              \end{itemize}
        \item Такой набор функций позволяет всесторонне оценить эффективность сравниваемых архитектур.
    \end{itemize}
\end{frame}

% Слайд 5: Результаты эксперимента и анализ
\begin{frame}{Результаты эксперимента и анализ}
    \begin{itemize}
        \item Для каждой архитектуры были подсчитаны средние значения метрик для разных временных интервалов:
              \begin{itemize}
                  \item \textbf{15 секунд:}
                        \begin{itemize}
                            \item GRUNet: MAPE = 6.78, \(R^2\) = 0.908, Память = 112.09 MB.
                            \item DenseNet: MAPE = 15.51, \(R^2\) = 0.844, Память = 36.92 MB.
                        \end{itemize}
                  \item \textbf{30 секунд:}
                        \begin{itemize}
                            \item GRUNet: MAPE = 3.90, \(R^2\) = 0.983, Память = 114.76 MB.
                            \item DenseNet: MAPE = 6.20, \(R^2\) = 0.969, Память = 39.75 MB.
                        \end{itemize}
                  \item \textbf{45 секунд:}
                        \begin{itemize}
                            \item GRUNet: MAPE = 2.81, \(R^2\) = 0.991, Память = 118.41 MB.
                            \item DenseNet: MAPE = 5.41, \(R^2\) = 0.978, Память = 44.89 MB.
                        \end{itemize}
              \end{itemize}
        \item Анализ показывает, что GRUNet обеспечивает более высокое качество аппроксимации при идентичных условиях обучения, несмотря на более высокие затраты памяти.
    \end{itemize}
\end{frame}

% Слайд 6: Заключение
\begin{frame}{Заключение}
    \begin{itemize}
        \item Внедрение архитектуры GRU в DEGANN значительно повышает точность аппроксимации ДУ.
        \item Преимущество GRUNet обусловлено её способностью учитывать последовательные зависимости, что приводит к более быстрой сходимости.
        \item Несколько более высокие затраты памяти не критичны для большинства практических задач.
    \end{itemize}
\end{frame}

\end{document}
